\begin{abstract}
  近年来，随着天地一体化信息网络、低轨卫星星座和在轨智能处理技术的快速发展，天基网络正由传统的数据传输网络向“通信、计算、存储与智能处理一体化”的分布式协同计算网络演进。卫星节点在灾害应急通信、遥感快速处理、空间态势感知、区域目标识别和边远区域信息服务等场景中承担着越来越多高时敏、高动态和高自治需求的计算任务。然而，受星载平台体积、功耗、能量循环、计算能力和星间链路条件限制，单星本地处理难以长期支撑高并发任务负载；传统星地集中处理模式又容易受到星地链路时延、可见窗口和回传带宽约束，难以满足任务快速闭环处理需求。因此，如何在高动态、异构和去中心化的天基环境下实现可信状态共享、多节点协同资源分配、并发感知任务卸载与平台化运行支撑，已成为天基协同计算网络亟需解决的关键问题。

  针对上述挑战，本文围绕天基网络协同计算卸载问题，提出一种融合区块链、多智能体强化学习、联邦学习与深度强化学习的双层智能管控与可信协同计算卸载方法体系。该体系以区块链作为分布式可信底座，在系统层实现多节点可信状态共享、协同博弈与资源分配，在任务层实现并发开销预测和实时卸载决策，并通过平台化架构将算法模型、状态同步、任务存证、模型治理和执行反馈统一纳入可运行、可验证、可审计的闭环环境中。本文主要研究内容与创新点如下：

  （1）提出基于区块链与多智能体强化学习的天基计算网络协同资源分配方法。针对去中心化天基环境中节点自利行为、信息不完全和资源竞争导致的负载失衡问题，将系统层协同资源分配建模为多智能体马尔可夫博弈，构建区块链增强的集中训练—分散执行多智能体强化学习框架。该方法以联盟链作为可信公共信息板，维护节点负载、链路状态、信誉评分和历史行为记录等公共摘要，并设计融合任务完成质量、系统公平性与信誉激励的混合奖励机制，引导异构卫星节点在分散自治条件下实现更稳定的协同资源分配。

  （2）提出基于联邦学习与深度强化学习的并发感知任务层计算卸载方法。针对多任务并发条件下 CPU 时间片、缓存、内存带宽和 I/O 队列竞争所引发的隐性执行代价，建立任务级并发开销模型，并构建“联邦学习预测 + DQN 卸载决策”的闭环机制。该方法利用聚类式联邦学习在不上传原始任务日志的前提下训练并发开销预测模型，通过联盟链维护模型版本和可信摘要，并将预测结果、本地状态和链上公共信息共同输入深度 Q 网络，使卫星终端能够在链上状态存在时滞和节点数据异构的条件下实现低开销、实时化、具备并发风险前瞻能力的任务卸载决策。

  （3）设计面向天基协同计算的区块链平台架构与运行机制。针对系统层协同优化和任务层实时决策缺乏统一承载环境的问题，构建面向动态星地网络的区块链协同计算平台。该平台通过轻量级联盟链共识、分层链结构、链上/链下协同机制和智能合约治理体系，实现关键状态可信同步、任务行为存证、模型版本溯源、节点信誉演化和执行反馈回流。平台将系统层全局协同资源分配方法与任务层并发感知卸载方法统一映射到工程化运行环境中，支撑天基协同计算从算法验证向可信部署和持续演化转化。

  （4）开展多场景仿真与系统级验证。围绕系统层协同资源分配、任务层并发感知卸载和区块链平台运行机制，设计收敛性、负载均衡、任务完成率、时延能耗、鲁棒性、消融分析和平台功能验证等实验。实验结果表明，本文所提方法能够在高动态链路、异构节点能力、任务非平稳到达和并发执行干扰等条件下有效降低平均任务时延和系统能耗，提高任务完成率、系统吞吐量与负载均衡水平，并增强分布式环境下的状态可信性、模型可追溯性和协同稳定性。

  综上，本文面向天基协同计算网络中“可信共享—系统协同—任务决策—平台承载”的完整链路，构建了融合区块链与人工智能的协同计算卸载方法体系。相关研究为提升受限资源条件下天基网络的任务处理效率、自治协同能力和可信运行水平提供了理论依据与技术支撑，也为未来构建高效、智能、可信的天基边缘计算网络奠定了基础。
\end{abstract}

% 如需手动控制换行连字符位置，可写 aa\-bb，详见
% https://bithesis.bitnp.net/faq/hyphen.html

\begin{abstractEn}
  In recent years, with the rapid development of space-ground integrated information networks, low Earth orbit satellite constellations, and on-orbit intelligent processing technologies, space-based networks are evolving from traditional data transmission networks into distributed collaborative computing networks integrating communication, computation, storage, and intelligent processing. Satellite nodes are undertaking increasingly time-sensitive, dynamic, and autonomous computing tasks in scenarios such as disaster emergency communication, rapid remote sensing processing, space situational awareness, regional target recognition, and information services for remote areas. However, constrained by onboard platform size, power consumption, energy cycles, computing capability, and inter-satellite link conditions, local processing on a single satellite can hardly sustain intensive concurrent workloads for a long time. Meanwhile, traditional space-ground centralized processing is limited by propagation latency, visibility windows, and feeder link bandwidth, making it difficult to satisfy the requirements of fast closed-loop task processing. Therefore, how to realize trusted state sharing, multi-node collaborative resource allocation, concurrency-aware task offloading, and platform-based operation support in highly dynamic, heterogeneous, and decentralized space environments has become a key challenge for space-based collaborative computing networks.

  To address these challenges, this thesis investigates collaborative computation offloading in space-based networks and proposes a dual-layer intelligent management and trusted collaborative computation offloading framework integrating blockchain, multi-agent reinforcement learning, federated learning, and deep reinforcement learning. In this framework, blockchain serves as a distributed trusted foundation. At the system layer, it supports trusted state sharing, collaborative game modeling, and multi-node resource allocation. At the task layer, it enables concurrency overhead prediction and real-time offloading decision-making. Furthermore, a platform-oriented architecture is designed to integrate algorithm models, state synchronization, task recording, model governance, and execution feedback into an operational, verifiable, and auditable closed-loop environment. The main research contents and contributions of this thesis are summarized as follows:

  (1) A blockchain and multi-agent reinforcement learning based collaborative resource allocation method for space-based computing networks is proposed. To address load imbalance caused by selfish node behaviors, incomplete information, and resource competition in decentralized space environments, the system-layer collaborative resource allocation problem is modeled as a multi-agent Markov game. A blockchain-enhanced centralized training and decentralized execution multi-agent reinforcement learning framework is constructed. In this method, the consortium blockchain acts as a trusted public information board to maintain public summaries such as node load, link status, reputation scores, and historical behavior records. A hybrid reward mechanism incorporating task completion quality, system fairness, and reputation incentives is further designed to guide heterogeneous satellite nodes toward stable collaborative resource allocation under decentralized autonomous execution.

  (2) A federated learning and deep reinforcement learning based concurrency-aware task-layer computation offloading method is proposed. To characterize the hidden execution costs caused by CPU time slicing, cache conflicts, memory bandwidth contention, and I/O queue competition under multi-task concurrency, a task-level concurrency overhead model is established, and a closed-loop mechanism of federated learning based prediction and DQN based offloading decision-making is designed. The proposed method employs clustered federated learning to train concurrency overhead prediction models without uploading raw task logs, maintains model versions and trusted summaries through blockchain, and feeds the predicted overhead, local state, and on-chain public information into a deep Q-network. As a result, satellite terminals can make low-overhead, real-time, and forward-looking offloading decisions under delayed on-chain states and heterogeneous node data distributions.

  (3) A blockchain platform architecture and operating mechanism for space-based collaborative computing are designed. To provide a unified execution environment for system-layer collaborative optimization and task-layer real-time decision-making, a blockchain-based collaborative computing platform for dynamic space-ground networks is constructed. Through lightweight consortium blockchain consensus, hierarchical chain structures, on-chain/off-chain collaboration, and smart contract based governance, the platform realizes trusted synchronization of key states, task behavior recording, model version traceability, node reputation evolution, and execution feedback. It maps the proposed system-layer global resource allocation method and task-layer concurrency-aware offloading method into an engineering-oriented operational environment, supporting the transition from algorithm validation to trusted deployment and continuous evolution.

  (4) Multi-scenario simulations and system-level verification are carried out. Experiments are designed for system-layer collaborative resource allocation, task-layer concurrency-aware offloading, and blockchain platform operation, including convergence analysis, load balancing, task completion rate, latency and energy consumption, robustness, ablation studies, and platform function verification. Experimental results show that the proposed methods can effectively reduce average task latency and system energy consumption, improve task completion rate, system throughput, and load balancing performance, and enhance state trustworthiness, model traceability, and collaborative stability in distributed environments with highly dynamic links, heterogeneous node capabilities, non-stationary task arrivals, and concurrent execution interference.

  In summary, this thesis constructs a blockchain and artificial intelligence integrated collaborative computation offloading framework for the complete chain of trusted sharing, system collaboration, task decision-making, and platform support in space-based collaborative computing networks. The research provides theoretical foundations and technical support for improving task processing efficiency, autonomous collaboration capability, and trusted operation under resource-constrained space environments, and lays a foundation for building efficient, intelligent, and trustworthy space-based edge computing networks in the future.
\end{abstractEn}